\subsection{The corrected cusp on the punctured bidisc}

\begin{proposition}[B.4.18; Proved; a partial cusp filling]
\label{push:B418}
Put
\[
 B^\times=(\Delta_q\times\Delta_r)\setminus\{(0,0)\}.
\]
There is a proper Hausdorff analytic family over
$B^\times$ with smooth total fourfold, flat over $r$,
whose restriction to $r\ne0$ is the actual prescribed
cusp family after $w=r^6$, and whose reduced
normal-crossings $r$-central fibre over $q\ne0$ has
component deck permutation the shift $-1$ on $\Z/12$
for the clockwise cusp generator $g_0$ of B.3.1.
Equivalently, the positive $q$ loop has shift $+1$.
In particular its $q$-fibre for every $r\ne0$ is the
actual reduced Mumford fibre, while its other axis has
the component transport required by B.3.2.

The overlap is an explicit holomorphic translation,
including the retained affine marking. This proposition
does not assert a proper completion over the missing
origin, nor the global finite-place overlap maps.
\end{proposition}

\begin{proof}
\emph{The fractional deck and its exact comparison.}
On the logarithmic cover of
$U=\Delta_q^*\times\Delta_r^*$ write
\[
 s=\frac{\log q}{2\pi i},\qquad
 t=\frac{\log r}{2\pi i},\qquad
 C(q)=\beta^{\rm part}(q)+s.
\]
The functions $\mu(q)$ and $C(q)$ are single-valued
and holomorphic at $q=0$. After the affine gauge of
B.4.9 the actual family is its Jacobian, with period
\[
 \boldsymbol\gamma(s,t)
       =(6\mu(q),C(q)-s-12t),\qquad
 \boldsymbol\delta=(0,1).
 \tag{B.4.18.1}
\]
The other periods are
$\mathbf u=(\tau,\mu)$ and $\mathbf w=(1,0)$.
Under $s\mapsto s+1$, they change by
$\boldsymbol\gamma\mapsto\boldsymbol\gamma-
\boldsymbol\delta$ and $\mathbf u\mapsto\mathbf u+\mathbf w$.
Under $t\mapsto t+1$, the period $\boldsymbol\gamma$
changes by $-12\boldsymbol\delta$.
These are integral changes of the period basis; the
ambient additive $\C^2$ coordinate is unchanged.

Set
\[
 b(s,t)=\frac1{12}\boldsymbol\gamma(s,t)
                         -\frac1{24}\boldsymbol\delta.
\]
The two exact identities
\[
 \sum_{k=0}^{11}b(s+k,t)
       =\boldsymbol\gamma(s,t)-6\boldsymbol\delta,
 \qquad
 b(s,t+1)-b(s,t)=-\boldsymbol\delta
 \tag{B.4.18.2}
\]
show that this gives a genuine order-twelve affine
$q$-deck action on the root cover $q=z^{12}$,
well defined around the $r$ loop. The first right-hand
side is an actual period. The scalar correction
$-\boldsymbol\delta/24$ is retained: the bare fraction
$\boldsymbol\gamma/12$ would not have the required
twelfth iterate.

There is an explicit holomorphic trivialization on $U$.
Put $L=s+12t$ and
\[
 F(s,t)=
 \left(\frac{\mu(q)L}{2},\,
       \frac{C(q)L}{12}-\frac{L^2}{24}+6t\right).
 \tag{B.4.18.3}
\]
Direct subtraction gives
\[
 F(s+1,t)-F(s,t)=b(s,t),\qquad
 F(s,t+1)-F(s,t)=\boldsymbol\gamma(s,t).
 \tag{B.4.18.4}
\]
The second difference is a period, not zero; that is
exactly what is needed for descent to the torus.
Thus $\zeta\mapsto\zeta+F$ is a holomorphic
isomorphism from the untwisted Jacobian to the
$+b$ affine torus torsor over $U$; its inverse
$\zeta\mapsto\zeta-F$ trivializes that torsor.
Replacing $b,F$ by $-b,-F$ gives the opposite
component rotation. If $\sigma$ denotes the actual
affine section of B.4.9, the map from the actual
family to the $+b$ torsor is explicitly
$\zeta\mapsto\zeta-\sigma+F$.
This retains the affine cusp marking, rather
than only identifying its fibrewise Jacobians.

\emph{A rotating compactification before changing the
$q$-boundary.}
We describe the root-cover model to verify the finite
action and properness. In cocharacter coordinates
$(y_1,y_2,h_z,h_r)$, its two period shears are
\[
 \begin{aligned}
 G_1 &: (y_1,y_2)\longmapsto(y_1+12h_z,y_2),\\
 G_2 &: (y_1,y_2)\longmapsto
                         (y_1,y_2+12h_z+12h_r).
 \end{aligned}
\]
The positive fractional deck has the shear
\[
 H:(y_1,y_2)\longmapsto(y_1,y_2+h_z+h_r),
 \qquad H^{12}=G_2,
 \tag{B.4.18.5}
\]
and rotates $z$ by $e^{2\pi i/12}$.
It commutes with $G_1$. Its unit factors are holomorphic
functions of $q$. For example, in the convention
$W=e^{-2\pi i\zeta_2}$, it multiplies $W$ by
\[
 r z\,e^{-2\pi i C(q)/12}e^{\pi i/12}
\]
and the first multiplicative coordinate by
$e^{\pi i\mu(q)}$. The twelve successive root rotations
contribute the factor
$e^{2\pi i(0+\cdots+11)/12}=-1$, and the scalar
correction contributes $e^{\pi i}=-1$.
They cancel. The twelfth iterate is consequently
the actual second period, including its unit, as
also follows additively from (B.4.18.2).

Use the unit-square cusp triangulation, with twelve
consecutive bottom cells in the $y_1$ direction and
the same bridge choice in every unit prism. Its
translates have bottom cell indices $(i,j)$ and top
interval $[j,j+1]$. Thus $H$ increases the bottom
second index and the top index by one, while $G_1$
shifts the first index by twelve. The fan is invariant
under both, and therefore under $G_2$.
For each bottom cell its three maximal cones are
the determinant-one cones of A.2.1. The usual
compact-core construction of A.1.2, now on the
$I_{12}$ elliptic root model and with twelve
plumbing intervals before the $G_2$ quotient,
gives a proper model on the root cover. Its
unfolded charts are precisely this fan. Thus
properness and proper discontinuity are supplied
by the compact core and collars, not inferred
from determinants alone.

The affine deck extends to this model because its
unit multiplication preserves every toric chart
and its shear preserves the fan. It is free.
On the $r$-central divisor a point lies on at most
two adjacent component labels in $\Z/12$; no
nonidentity translation can preserve that set.
This covers their intersection over the common
base origin as well. On the $q$-central divisor
with $r\ne0$, the bottom triangle has second
component indices contained in $\{j,j+1\}$,
again not preserved by a nonzero translation
modulo twelve. Off both axes the root-base action
is free. Its finite quotient is therefore smooth,
proper over $(q,r)$, flat over $r$, and has a
reduced normal-crossings $r$-central divisor.

For completeness, for the negative rotation keep the same positive
root-base rotation and use the inverse translation
$-b$. The resulting valuation shear is $H^{-1}$;
its twelfth iterate is $G_2^{-1}$. It preserves
the same fan and has the same freeness properties.
This is the opposite convention and is not the model
used below. In the retained convention the clockwise
cusp generator is $g_0:s\mapsto s-1$. Thus the
positive fractional deck $b$ has shift $+1$ for
$s\mapsto s+1$ and shift $-1$ for $g_0$, exactly
as required by B.3.2. The finite shifts $4,-3$
and this cusp shift all refer to the same clockwise
deck-identification convention.

Denote this rotating model by $\mathcal R$.
Its $q$-fibre for $r\ne0$ has multiplicity twelve.
Indeed the original $z$-central divisor is reduced
and the finite action has trivial generic isotropy
there, while $q=z^{12}$. This is exactly the
limitation established in B.3.3. We do not identify
this full $q$-boundary with the prescribed cusp.

\emph{The actual two-open gluing.}
On the base open set $q\ne0$ use $\mathcal R$ with
the positive fractional deck $b$. On the base open set $r\ne0$
use the actual reduced Mumford cusp family.
Their common base is $U$, where (B.4.18.3)--
(B.4.18.4), composed with B.4.9, give the specified
holomorphic isomorphism. These two open sets cover
$B^\times$, and there is no additional triple
overlap. Glue by that isomorphism.

The result is Hausdorff: points over one base
point are separated in one of the two given
Hausdorff models, and points with different
base images are separated in the base.
Properness is local on the target, so follows
from the properness of those two restrictions.
Smoothness and flatness over $r$ are also local.
The restriction over $r\ne0$ is, by construction,
the actual prescribed cusp including its reduced
$q$-fibre. Over $r=0,q\ne0$, the construction has
the required twelve-cycle and rotation $-1$.
This proves every asserted property.

The exact symbolic differences, finite norm,
scalar phase, commuting shears and unit-prism
determinants are saved in
\texttt{code/ab\_push\_corrected\_cusp\_overlap.py}
and its output. No completion at the origin
follows from these finite calculations.
\end{proof}

\begin{remark}
The next local construction must complete this
proper family over the single missing base point.
The comparison (B.4.18.3) contains quadratic
logarithms. Neither a meromorphic graph closure
nor a common toroidal modification is thereby
available. For fixed $r\ne0$, the multiple
$q$-boundary of $\mathcal R$ cannot be substituted
for the actual reduced boundary: the valuation
argument of B.3.7 excludes a proper birational
comparison over $q$ between them. A successful
origin cap must retain the actual two-open
comparison just constructed and prove its own
analytic extension and properness.
\end{remark}
